\section{Experiments}
\label{sec:experiments}

\subsection{Experimental Setup}
\label{sec:exp_setup}

We evaluate vision token pruning across six VLA models spanning three connector families on the LIBERO benchmark~\cite{liu2023libero}. LIBERO provides four task suites of increasing complexity: Object (object manipulation), Spatial (spatial reasoning), Goal (goal-conditioned tasks), and Long (long-horizon multi-step tasks). Each suite contains 10 tasks evaluated over 50 episodes, yielding 500 episodes per condition.

Four models are evaluated in closed-loop robotic control: $\pi_0.5$~\cite{black2025pi05} (4B parameters, linear projection with $\times\!\sqrt{d}$ scaling, 768 tokens from three cameras, CV$\,{=}\,$0.055), X-VLA~\cite{zheng2025xvla} (880M, linear projection followed by a 12-layer BART encoder~\cite{lewis2020bart}, 50 tokens, CV$\,{=}\,$0.151), OpenVLA~\cite{kim2024openvla} (7B, three-layer MLP projector, 256 tokens, CV$\,{=}\,$0.215), and SmolVLA~\cite{shukor2025smolvla} (450M, pixel-shuffle $4{\times}4$ folding, 128 tokens, CV$\,{=}\,$0.314). Two additional models are analyzed via offline metrics only: $\pi_0$~\cite{black2024pi0} (3.3B, linear projection, 256 tokens, CV$\,{=}\,$0.038) and OFT~\cite{kim2025openoft} (7B, MLP projector, 512 tokens, CV$\,{=}\,$0.219). Table~\ref{tab:model_properties} summarizes model properties and connector details.

The pruning protocol applies post-connector token pruning at retention ratios $k \in \{0.3, 0.5, 0.7, 0.85, 0.9, 0.95\}$. Two strategies bracket the robustness range: \textbf{L2 norm selection} ranks tokens by L2 norm and retains the top-$k$ fraction, providing an importance-based favorable test; \textbf{random subsampling} selects tokens uniformly at random, providing a conservative lower bound. Both strategies physically remove tokens and re-index position IDs. All results report mean success rate (SR) over 500 episodes with Wilson 95\% confidence intervals for single-seed conditions.

\subsection{Two-Axis Framework Overview}
\label{sec:two_axis}

The pruning response of a VLA model depends on two axes: the degree of token independence imposed by the vision connector, and the complexity of the downstream manipulation task. Models with token-independent connectors (MLP projectors, linear projections without deep post-processing) degrade gracefully under pruning, while models with token-entangling connectors exhibit sharp, task-dependent collapse. Figure~\ref{fig:teaser} illustrates this two-axis structure.

Table~\ref{tab:main_results} presents closed-loop success rates across four models, four LIBERO suites, two pruning methods, and four retention ratios. The data reveal four qualitatively distinct response patterns spanning the connector spectrum (Figure~\ref{fig:pruning_overview}).

\textbf{$\pi_0.5$ (linear, CV$\,{=}\,$0.055).} On Object, L2 norm selection yields 98.2\% at $k{=}0.95$, 98.2\% at $k{=}0.9$, and 98.0\% at $k{=}0.85$ (a plateau within 0.2\,pp of the 98.2\% baseline), then drops sharply to 55.0\% at $k{=}0.7$. Random selection follows a different trajectory: 69.4\% at $k{=}0.95$, 10.8\% at $k{=}0.9$, 75.5\% at $k{=}0.7$, and 68.9\% at $k{=}0.5$. The low CV implies near-uniform token norms, and the method choice produces gaps as large as 87.4\,pp (at $k{=}0.9$). On Spatial, $\pi_0.5$ maintains 97.8\% at $k{=}0.95$ (L2) and degrades to 56.9\% at $k{=}0.5$ (L2) versus 71.6\% (random). The pattern is consistent: low-CV models are not pruning-immune but uniformly critical, where the choice of which tokens to remove matters more than the number removed.

\textbf{X-VLA (linear + BART encoder, CV$\,{=}\,$0.151).} X-VLA exhibits threshold-graded behavior with a pronounced interaction between retention ratio and task complexity. On Object, L2 norm selection preserves 99.0\% at $k{=}0.95$ and 97.6\% at $k{=}0.9$, then falls sharply to 68.0\% at $k{=}0.85$ with non-monotonic recovery to 89.0\% at $k{=}0.8$ before continued decline to 56.2\% at $k{=}0.7$. The same model at the same $k{=}0.9$ achieves only 70.2\% on Goal, a 27.8\,pp difference driven entirely by task complexity. On Long, removing just 3 of 50 tokens ($k{=}0.95$) collapses SR from 94.0\% to 47.0\%. The 12-layer BART encoder creates cross-token dependencies that amplify under compositional task demands.

\textbf{OpenVLA (MLP, CV$\,{=}\,$0.215).} OpenVLA degrades smoothly across all retention ratios on Object: 70.6\% (L2) and 69.2\% (random) at $k{=}0.95$, 62.0\% and 63.4\% at $k{=}0.9$, 49.0\% and 48.5\% at $k{=}0.7$, 37.4\% and 35.7\% at $k{=}0.5$. The near-zero method gap (within 2\,pp at all retention ratios) reflects the MLP projector's per-patch independence: with a genuine importance gradient (CV$\,{=}\,$0.215), L2 norm ranking and random selection yield comparable spatial coverage. Degradation is proportional to the fraction removed, with no threshold effects.

\textbf{SmolVLA (pixel-shuffle, CV$\,{=}\,$0.314).} SmolVLA collapses to 0.0\% SR at every retention ratio across all four suites and both pruning methods, including $k{=}0.95$ where only 6 of 128 tokens are removed. This catastrophic cliff is detailed in Section~\ref{sec:pixel_shuffle}.

\begin{table*}[t]
\centering
\small
\caption{Closed-loop success rates (\%) on LIBERO under L2 norm (L2) and random (R) token pruning at four retention ratios. Best result per row in \textbf{bold}. SmolVLA collapses to 0.0\% across all conditions. ``--'' indicates data not yet collected.}
\label{tab:main_results}
\begin{tabular}{llccccc}
\toprule
Model & Suite & Baseline & $k{=}0.95$ (L2 / R) & $k{=}0.9$ (L2 / R) & $k{=}0.7$ (L2 / R) & $k{=}0.5$ (L2 / R) \\
\midrule
\multirow{4}{*}{$\pi_0.5$}
& Object  & \textbf{98.2} & 98.2 / 69.4 & \textbf{98.2} / 10.8 & 55.0 / 75.5 & 33.4 / 68.9 \\
& Spatial & \textbf{97.2} & \textbf{97.8} / 76.6 & 95.0 / 54.1 & 67.7 / 74.1 & 56.9 / 71.6 \\
& Goal    & \textbf{96.4} & 89.4 / -- & 90.0 / -- & -- / 80.6 & 59.2 / 74.4 \\
& Long    & \textbf{91.6} & -- / 61.2 & 81.0 / 31.8 & 60.0 / 64.6 & 50.2 / 50.2 \\
\midrule
\multirow{4}{*}{X-VLA}
& Object  & \textbf{99.2} & 99.0 / 98.8 & 97.6 / 97.8 & 56.2 / 65.2 & 42.6 / 70.6 \\
& Spatial & \textbf{91.6} & 71.6 / -- & 78.8 / -- & 72.0 / -- & 58.8 / 71.4 \\
& Goal    & \textbf{98.0} & 53.8 / -- & 70.2 / -- & 53.4 / -- & 61.0 / -- \\
& Long    & \textbf{94.0} & 47.0 / -- & 27.5 / -- & -- / -- & -- / -- \\
\midrule
OpenVLA & Object & \textbf{74.0} & 70.6 / 69.2 & 62.0 / 63.4 & 49.0 / 48.5 & 37.4 / 35.7 \\
\midrule
\multirow{4}{*}{SmolVLA}
& Object  & \textbf{59.2} & 0.0 / 0.0 & 0.0 / 0.0 & 0.0 / 0.0 & 0.0 / 0.0 \\
& Spatial & \textbf{44.6} & 0.0 / 0.0 & 0.0 / 0.0 & 0.0 / 0.0 & 0.0 / 0.0 \\
& Goal    & \textbf{60.8} & 0.0 / 0.0 & 0.0 / 0.0 & 0.0 / 0.0 & 0.0 / 0.0 \\
& Long    & \textbf{36.6} & 0.0 / 0.0 & 0.0 / 0.0 & 0.0 / 0.0 & 0.0 / 0.0 \\
\bottomrule
\end{tabular}
\end{table*}

\subsection{Pixel-Shuffle Entanglement}
\label{sec:pixel_shuffle}

SmolVLA's pixel-shuffle connector applies $4{\times}4$ spatial folding, concatenating 16 contiguous patches into the channel dimension of a single token. This operation creates irreversible cross-token spatial dependencies: each token encodes a $4{\times}4$ receptive field, and removing any token eliminates information from all 16 constituent patches with no redundancy in the remaining tokens.

The consequence is catastrophic. SmolVLA collapses to 0.0\% SR at $k{=}0.95$ (removing 6 of 128 tokens, a 5\% reduction) across all four LIBERO suites (Object, Spatial, Goal, Long) and both pruning methods (L2 norm and random). We verified this result on two independent checkpoints: our LoRA fine-tuned reproduction (44.8\% SR baseline) and the official SmolVLA checkpoint (59.2\% SR baseline). Both exhibit identical collapse. The cliff is architectural, not a training artifact.

The pixel-shuffle mechanism explains why the collapse is method-independent. Because all tokens carry spatially entangled information from non-overlapping $4{\times}4$ grids, no scoring function can identify ``dispensable'' tokens. L2 norm ranking and random selection fail equally because the problem is not which tokens are removed but that any removal breaks the spatial completeness the downstream LLM requires. In contrast, MLP projectors (OpenVLA) and linear projections ($\pi_0.5$) map patches independently. Each token retains per-patch spatial locality, so partial removal does not destroy information from unrelated spatial positions.

\subsection{When Does Pruning Method Matter? CV as Sensitivity Indicator}
\label{sec:cv_sensitivity}

Token importance CV (coefficient of variation of L2 norms) correlates with the retention ratio at which the choice of pruning method becomes critical. At mild pruning ($k{=}0.95$), a clear pattern emerges across the three closed-loop models with distinct CV values.

For $\pi_0.5$ (CV$\,{=}\,$0.055), the method gap at $k{=}0.95$ is substantial: 98.2\% (L2) versus 69.4\% (random) on Object, a 28.8\,pp difference. On Spatial, the gap is 21.4\,pp (97.8\% vs.\ 76.6\%). The near-uniform token norms cause L2 norm ranking to exploit small spatial correlations rather than genuine importance differences, while random selection preserves the uniform spatial coverage these models depend on.

For X-VLA (CV$\,{=}\,$0.151), the method gap at $k{=}0.95$ is minimal on Object (0.2\,pp; Table~\ref{tab:main_results}). Multi-seed evaluation across suites yields an average gap of 1.2\,pp. The BART encoder produces tokens with a moderate importance gradient that both methods preserve at this mild pruning level.

For OpenVLA (CV$\,{=}\,$0.215), the method gap at $k{=}0.95$ is 1.4\,pp on Object (70.6\% vs.\ 69.2\%). The MLP projector creates a clear importance hierarchy that renders both selection strategies nearly equivalent.

The contrast is stark: $\pi_0.5$ exhibits a 28.8\,pp method gap where X-VLA and OpenVLA show gaps below 1.5\,pp. Low-CV models are an order of magnitude more sensitive to pruning method choice than higher-CV models at mild retention ratios. Figure~\ref{fig:method_gap} visualizes this relationship.

At aggressive compression ($k{=}0.5$), the method-independence observed in moderate-CV models breaks down. On Object, the gap $\Delta$(R$-$L2) scales monotonically with CV: $\pi_0.5$ shows $+$35.5\,pp (random: 68.9\%, L2: 33.4\%), X-VLA shows $+$28.0\,pp (random: 70.6\%, L2: 42.6\%), and OpenVLA shows $-$1.7\,pp (random: 35.7\%, L2: 37.4\%). At 50\% retention, random selection outperforms L2 norm selection for both low-CV and moderate-CV models because random sampling preserves spatial diversity, while L2 norm ranking selects correlated high-norm tokens that cluster in similar image regions. Only OpenVLA, with CV above 0.2, produces a genuine importance gradient that L2 can exploit without spatial clustering artifacts.

\subsection{Cross-Suite Task Complexity Effects}
\label{sec:cross_suite}

LIBERO's four suites form a complexity gradient: Object and Spatial require single-step manipulation with explicit spatial cues, while Goal and Long demand multi-step planning, goal inference, or extended temporal coordination. This gradient interacts with pruning tolerance in a connector-dependent manner.

$\pi_0.5$ shows moderate complexity amplification under L2 norm pruning at $k{=}0.95$: Object retains 98.2\% (matching the baseline), Spatial retains 97.8\% (within noise of the 97.2\% baseline), and Goal drops to 89.4\% ($-$7.0\,pp from 96.4\%). Simple tasks absorb mild pruning with negligible loss; compositional tasks expose measurable degradation even at 5\% token removal.

X-VLA exhibits dramatic complexity amplification. At $k{=}0.9$ with L2 norm selection, Object retains 97.6\% ($-$1.6\,pp from the 99.2\% baseline), while Goal drops to 70.2\% ($-$27.8\,pp). The degradation factor between suites is 17$\times$ (27.8 / 1.6), driven entirely by task complexity with the same model, same retention ratio, and same pruning method. The effect is even more pronounced at $k{=}0.95$ on Long: removing just 3 of 50 tokens collapses SR from 94.0\% to 47.0\%, a 47.0\,pp loss, while Object at the same ratio loses only 0.2\,pp.

Token-entangling connectors (X-VLA's BART encoder) amplify task complexity because compositional tasks require integrating information across spatial locations. When the BART encoder creates cross-token dependencies, removing even a few tokens disrupts the distributed representations that multi-step tasks rely on. Token-independent connectors ($\pi_0.5$'s linear projection, OpenVLA's MLP) show moderate complexity amplification because each token's contribution is self-contained: removing tokens reduces spatial coverage proportionally, and simple tasks tolerate the resulting gaps.

\subsection{VTR Recovery}
\label{sec:vtr_recovery}

Vision Token Re-encoder (VTR) addresses the catastrophic cliff through learned re-encoding rather than hard deletion. VTR inserts a lightweight cross-attention module ($\sim$22M parameters, approximately 5\% of SmolVLA) between the pixel-shuffle connector and the LLM backbone. $M$ learned queries attend over the full set of 128 source tokens and produce $M$ compressed output representations, where $M \in \{16, 32, 64, 96, 128\}$. The $M{=}128$ configuration serves as a passthrough control that re-encodes without compressing.

Table~\ref{tab:vtr_results} reports VTR results on SmolVLA for LIBERO-Object (3 seeds, 500 episodes each). Cross-attention VTR recovers substantial performance where hard pruning yields 0.0\%: $M{=}16$ (87.5\% compression) achieves 56.9 $\pm$ 1.3\%, $M{=}32$ (75\%) achieves 60.1 $\pm$ 1.3\%, $M{=}64$ (50\%) achieves 62.0 $\pm$ 2.2\%, $M{=}96$ (25\%) achieves 54.5 $\pm$ 3.2\%, and $M{=}128$ (0\%, passthrough) achieves 58.0 $\pm$ 3.5\%. The SmolVLA-Base (LoRA fine-tuned) baseline without VTR is 44.8\%, and the official SmolVLA baseline is 59.2\%.

\begin{table}[t]
\centering
\small
\caption{VTR re-encoding on SmolVLA (LIBERO-Object). Cross-attention and MLP adapter results report mean $\pm$ std over 3 seeds (500 episodes each). The M-curve is flat: $M{=}16$ (87.5\% compression) matches $M{=}128$ (passthrough) within 1.1\,pp. The baseline control (100K extra training, no compressor) confirms the gain is architectural. ``Compression'' = fraction of tokens removed.}
\label{tab:vtr_results}
\begin{tabular}{lccr}
\toprule
Configuration & $M$ & Compression & SR (\%) \\
\midrule
\textit{SmolVLA-Base (LoRA)} & 128 & 0\% & 44.8 \\
\textit{SmolVLA official} & 128 & 0\% & 59.2 \\
\textit{Baseline control (+100K steps)} & 128 & 0\% & 45.0 \\
\textit{Hard pruning (any $k$)} & -- & 5--50\% & 0.0 \\
\midrule
Cross-attn VTR & 16 & 87.5\% & 56.9 $\pm$ 1.3 \\
Cross-attn VTR & 32 & 75\% & 60.1 $\pm$ 1.3 \\
Cross-attn VTR & 64 & 50\% & \textbf{62.0 $\pm$ 2.2} \\
Cross-attn VTR & 96 & 25\% & 54.5 $\pm$ 3.2 \\
Cross-attn VTR & 128 & 0\% & 58.0 $\pm$ 3.5 \\
\midrule
MLP adapter ($\sim$22M) & 32 & 75\% & 63.4 $\pm$ 2.6 \\
\bottomrule
\end{tabular}
\end{table}

The M-curve is remarkably flat: performance at $M{=}16$ (56.9\%) falls within 1.1\,pp of $M{=}128$ (58.0\%), and all five $M$ values cluster within a 7.5\,pp range. This flatness confirms that the operative factor is learned re-encoding of spatially entangled tokens, not the compression ratio itself. The VTR module transforms pixel-shuffle representations into a format the downstream LLM can process robustly, and this transformation is effective regardless of how many output tokens carry the result.

A parameter-matched MLP adapter ($\sim$22M parameters) achieves 63.4 $\pm$ 2.6\% at $M{=}32$, comparable to cross-attention VTR ($p = 0.15$). This confirms that re-encoding itself, rather than the specific attention mechanism, is the key architectural intervention. The baseline control (100K additional training steps without the compressor module) achieves only 45.0\% SR, within 0.2\,pp of the 44.8\% baseline. The performance gain is therefore attributable to the VTR architecture, not to additional optimization.

Cross-suite evaluation at $M{=}32$ confirms that VTR recovery generalizes beyond Object: Spatial achieves 53.2\% and Goal achieves 56.4\%, both recovered from 0.0\% under hard pruning. VTR provides a targeted remedy for pixel-shuffle entanglement. By breaking the spatial dependencies that make hard deletion catastrophic, VTR enables both re-encoding recovery and downstream token reduction for deployment efficiency.
